% Generated draft: requires mathematical and source review.
\documentclass[11pt]{article}
\usepackage[margin=25mm]{geometry}
\usepackage{fontspec}
\setmainfont{Latin Modern Roman}
\usepackage{amsmath,amssymb,mathtools,unicode-math}
\setmathfont{Latin Modern Math}
\usepackage{enumitem,needspace,xurl,hyperref}
\setlength{\parindent}{0pt}
\setlength{\parskip}{5pt}
\setlength{\emergencystretch}{3em}
\newcommand{\N}{\mathbb{N}}
\newcommand{\Z}{\mathbb{Z}}
\newcommand{\Q}{\mathbb{Q}}
\newcommand{\R}{\mathbb{R}}
\newcommand{\C}{\mathbb{C}}
\newcommand{\F}{\mathbb{F}}
\newcommand{\A}{\mathbb{A}}
\newcommand{\PP}{\mathbb{P}}
\newcommand{\E}{\mathbb{E}}
\newcommand{\eps}{\varepsilon}
\newcommand{\dd}{\,\mathrm{d}}
\begin{document}
\section*{UnsolvedMath ID: 6700084}
\textbf{Source problem code:} catalogue entry
\subsection{Definitions and mathematical statement}
Let \(X\) be an \(n\)-dimensional Alexandrov space with curvature bounded below by \(-1\) in the lower triangle-comparison sense. This is not the CAT\((-1)\) upper-curvature condition. Assume the source's volumic scalar lower bound \(\operatorname{Sc}^{\mathrm{vol}}_n(X)\ge0\). For a continuous map \(\Phi:X\to T^n\), the source explicitly defines nonzero degree by nonvanishing of the induced top-homology map. With ordinary integral coefficients, this is
\[
\Phi_*:H_n(X;\Z)\longrightarrow H_n(T^n;\Z)\cong\Z,
\qquad\Phi_*\ne0.
\]
No Lipschitz restriction on \(\Phi\) or distinguished fundamental class for the singular space is added.

The conjecture is
\[
\operatorname{curv}_{\mathrm{Alex}}(X)\ge-1,
\quad\operatorname{Sc}^{\mathrm{vol}}_n(X)\ge0,
\quad\exists\Phi\text{ as above}
\quad\Longrightarrow\quad
\widetilde X\cong\R^n\text{ isometrically},
\]
where the universal cover carries the lifted length metric and \(\R^n\) its Euclidean metric. The volumic predicate and homology coefficient convention are not defined in the extraction. Compactness and boundary assumptions are also not stated; they are not silently imposed merely because a nonzero top-homology map is required. The source itself says the volumic notion is provisional and intended to inspire a more satisfactory scalar-curvature definition. The supplied summary reports no located resolution of this singular-space extension of torus rigidity.

Adjacent section 27 begins Fredholm coarea and stable K-area. It allows complex Hilbert bundles with unitary connections described through parallel transport. Its proposed curvature control is that holonomy rotation angles around suitable surfaces satisfy
\[
|\alpha|\le\int_S C(s)\,\mathrm dA(s).
\]
For a K-cohomology class \(\kappa\) represented by a Fredholm bundle map \(F:L_1\to L_2\), its coarea is described as the infimum of constants \(c\) admitting such a representation and connections on both bundles with curvature norms at most \(c\). In the compact-support version, \(F\) must be a unitary isomorphism and preserve the connections outside a compact set. The source then proposes dual Fredholm K-area on K-homology, ordinary homology, and infinite-support homology, but ends before the formula. The admissible surfaces, infinite-dimensional curvature convention, and missing duality formula cannot be completed from the record. This adjacent framework is not an extra hypothesis of the Alexandrov conjecture.
\subsection{Short English statement}
Does a nonnegative-volumic-scalar-curvature Alexandrov space mapping nontrivially in top homology to the torus have Euclidean universal cover? Preserve the lower-curvature sign and the explicitly homological degree condition; the adjacent Fredholm framework is incomplete.
\subsection{Sources}
\begin{itemize}
\item Local UnsolvedMath registry, permanent ID \(6700084\), exact source code \texttt{catalogue entry}; printed scalar-curvature tag [?86], occurrence 87, ``Geroch for Alexandrov Spaces.'' Proposer Misha Gromov, 2017.
\item Gromov, \emph{101 Questions, Problems and Conjectures around Scalar Curvature} (2017), recorded PDF page 84; \url{https://www.ihes.fr/~gromov/wp-content/uploads/2018/08/101-problemsOct1-2017.pdf}. Recorded access 2026-07-29; pdf-text-normalized extraction.
\item Adjacent section 27 refers back to sections 9, 12 and 13 for curvature and K-area conventions, absent here; its final duality formula is truncated. The supplied OPEN-TRIAGE research summary gives no additional precise resolution citation.
\end{itemize}
\end{document}
